\documentclass[aspectratio=169]{beamer}
\usetheme{Madrid}
\usecolortheme{default}
\usepackage[utf8]{inputenc}
\usepackage[spanish,es-minimal]{babel}
\usepackage{amsmath}
\usepackage{amssymb}
\usepackage{booktabs}
\usepackage{multirow}
\usepackage{graphicx}

\definecolor{unsblue}{RGB}{30,60,120}
\setbeamercolor{structure}{fg=unsblue}
\setbeamercolor{palette primary}{bg=unsblue,fg=white}
\setbeamercolor{palette secondary}{bg=unsblue!80,fg=white}
\setbeamercolor{palette tertiary}{bg=unsblue!60,fg=white}

\setbeamertemplate{navigation symbols}{}
\setbeamertemplate{footline}[frame number]

\title[MAD1 - Clase 23]{Métodos de Análisis de Datos 1}
\subtitle{Clase 23: Implementación Computacional de ANOVA y Diseños Experimentales}
\author{Prof. José}
\institute{Departamento de Matemática \\ Universidad Nacional del Sur}
\date{}

\begin{document}

\begin{frame}
\titlepage
\end{frame}

\begin{frame}{Contenidos de la Clase}
\tableofcontents
\end{frame}

\section{Introducción}

\begin{frame}{Objetivos de la Clase}
En esta clase aprenderemos a implementar en R y Python:

\begin{enumerate}
\item \textbf{ANOVA de dos factores:}
\begin{itemize}
\item Con y sin interacción
\item Gráficos de interacción
\item Comparaciones post-hoc
\end{itemize}

\vspace{0.2cm}

\item \textbf{Diseños experimentales clásicos:}
\begin{itemize}
\item Diseño Completamente Aleatorizado (DCA)
\item Diseño en Bloques Completos Aleatorizados (DBCA)
\item Diseño en Cuadrado Latino (DCL)
\end{itemize}

\vspace{0.2cm}

\item \textbf{Verificación de supuestos:}
\begin{itemize}
\item Normalidad, homocedasticidad, independencia
\end{itemize}
\end{enumerate}
\end{frame}

\begin{frame}{Paquetes y Librerías}
\textbf{En Python:}
\begin{itemize}
\item \texttt{pandas}: manipulación de datos
\item \texttt{numpy}: cálculos numéricos
\item \texttt{scipy.stats}: pruebas estadísticas
\item \texttt{statsmodels}: modelos estadísticos (ANOVA)
\item \texttt{matplotlib, seaborn}: visualización
\item \texttt{pingouin}: estadística (alternativa a statsmodels)
\end{itemize}

\vspace{0.3cm}

\textbf{En R:}
\begin{itemize}
\item Base R: \texttt{aov()}, \texttt{lm()}, \texttt{anova()}
\item \texttt{car}: pruebas adicionales (Levene, Type II/III SS)
\item \texttt{agricolae}: comparaciones múltiples, diseños
\item \texttt{emmeans}: medias marginales estimadas
\item \texttt{ggplot2}: visualización
\end{itemize}
\end{frame}

\section{ANOVA de Dos Factores}

\subsection{Ejemplo: Rendimiento de Cultivos}

\begin{frame}[fragile]{Datos del Ejemplo}
\small
Rendimiento de cultivos (kg/ha) según:
\begin{itemize}
\item Factor A (Fertilizante): F1, F2, F3
\item Factor B (Riego): Bajo, Alto
\item 4 réplicas por combinación
\end{itemize}

\vspace{0.2cm}

\textbf{Estructura de datos:}
\begin{verbatim}
   Fertilizante  Riego  Rendimiento
1            F1   Bajo           45
2            F1   Bajo           47
3            F1   Bajo           46
...
24           F3   Alto           65
\end{verbatim}
\end{frame}

\begin{frame}[fragile]{Python: Cargar y Preparar Datos}
\small
\begin{verbatim}
import pandas as pd
import numpy as np
from scipy import stats
import statsmodels.api as sm
from statsmodels.formula.api import ols

# Crear los datos
data = {
    'Fertilizante': ['F1']*8 + ['F2']*8 + ['F3']*8,
    'Riego': ['Bajo']*4 + ['Alto']*4 +
             ['Bajo']*4 + ['Alto']*4 +
             ['Bajo']*4 + ['Alto']*4,
    'Rendimiento': [45,47,46,48, 55,57,56,58,
                    52,54,53,55, 58,60,59,61,
                    48,50,49,51, 62,64,63,65]
}
df = pd.DataFrame(data)
\end{verbatim}
\end{frame}

\begin{frame}[fragile]{Python: ANOVA de Dos Factores}
\small
\begin{verbatim}
# Modelo con interacción
modelo = ols('Rendimiento ~ C(Fertilizante) + C(Riego) 
              + C(Fertilizante):C(Riego)', 
              data=df).fit()

# Tabla ANOVA
tabla_anova = sm.stats.anova_lm(modelo, typ=2)
print(tabla_anova)
\end{verbatim}

\vspace{0.2cm}

\textbf{Salida:}
\begin{verbatim}
                              sum_sq    df         F    PR(>F)
C(Fertilizante)            126.083     2    18.953  0.000049
C(Riego)                   600.000     1   180.339  <0.000001
C(Fertilizante):C(Riego)    84.083     2    12.631  0.000293
Residual                    59.833    18
\end{verbatim}
\end{frame}

\begin{frame}[fragile]{R: ANOVA de Dos Factores}
\small
\begin{verbatim}
# Crear los datos
rendimiento <- c(45,47,46,48, 55,57,56,58,
                 52,54,53,55, 58,60,59,61,
                 48,50,49,51, 62,64,63,65)
fertilizante <- factor(rep(c("F1","F2","F3"), each=8))
riego <- factor(rep(rep(c("Bajo","Alto"), each=4), 3))

# ANOVA con interacción
modelo <- aov(rendimiento ~ fertilizante * riego)
summary(modelo)
\end{verbatim}

\vspace{0.2cm}

\textbf{Salida:}
\begin{verbatim}
                      Df Sum Sq Mean Sq F value   Pr(>F)    
fertilizante           2  126.1   63.04  18.953 4.93e-05 ***
riego                  1  600.0  600.00 180.339  < 2e-16 ***
fertilizante:riego     2   84.1   42.04  12.631 0.000293 ***
Residuals             18   59.8    3.32                     
\end{verbatim}
\end{frame}

\begin{frame}{Interpretación de Resultados}
\textbf{Conclusiones del ejemplo:}

\begin{enumerate}
\item \textbf{Interacción significativa} ($p < 0.001$):
\begin{itemize}
\item El efecto del fertilizante depende del nivel de riego
\item No interpretar efectos principales por separado
\end{itemize}

\vspace{0.2cm}

\item \textbf{Efecto del riego} ($p < 0.001$):
\begin{itemize}
\item Riego alto produce mayor rendimiento
\end{itemize}

\vspace{0.2cm}

\item \textbf{Efecto del fertilizante} ($p < 0.001$):
\begin{itemize}
\item Hay diferencias entre fertilizantes
\item Pero el efecto depende del riego (interacción)
\end{itemize}
\end{enumerate}

\vspace{0.2cm}
\textbf{Próximo paso:} Graficar la interacción y realizar comparaciones post-hoc.
\end{frame}

\subsection{Gráficos de Interacción}

\begin{frame}[fragile]{Python: Gráfico de Interacción}
\small
\begin{verbatim}
import matplotlib.pyplot as plt
import seaborn as sns

# Calcular medias por combinación
medias = df.groupby(['Fertilizante', 'Riego'])['Rendimiento'].mean()
medias = medias.reset_index()

# Gráfico de interacción
fig, ax = plt.subplots(figsize=(10, 6))
for fert in ['F1', 'F2', 'F3']:
    datos_fert = medias[medias['Fertilizante'] == fert]
    ax.plot(datos_fert['Riego'], datos_fert['Rendimiento'],
            marker='o', label=fert, linewidth=2, markersize=8)

ax.set_xlabel('Nivel de Riego', fontsize=12)
ax.set_ylabel('Rendimiento Medio (kg/ha)', fontsize=12)
ax.set_title('Gráfico de Interacción', fontsize=14)
ax.legend(title='Fertilizante')
ax.grid(True, alpha=0.3)
plt.tight_layout()
plt.show()
\end{verbatim}
\end{frame}

\begin{frame}[fragile]{R: Gráfico de Interacción}
\small
\begin{verbatim}
# Gráfico de interacción con ggplot2
library(ggplot2)

datos <- data.frame(rendimiento, fertilizante, riego)
medias <- aggregate(rendimiento ~ fertilizante + riego,
                    data = datos, FUN = mean)

ggplot(medias, aes(x = riego, y = rendimiento,
                   color = fertilizante, group = fertilizante)) +
  geom_line(linewidth = 1.2) +
  geom_point(size = 4) +
  labs(x = "Nivel de Riego",
       y = "Rendimiento Medio (kg/ha)",
       title = "Gráfico de Interacción",
       color = "Fertilizante") +
  theme_minimal() +
  theme(text = element_text(size = 12))
\end{verbatim}
\end{frame}

\begin{frame}{Interpretación del Gráfico}
\textbf{Observaciones del gráfico de interacción:}

\begin{itemize}
\item \textbf{Líneas no paralelas}: confirma interacción significativa
\item Con riego bajo: F2 es mejor que F1 y F3
\item Con riego alto: F3 es el mejor, F2 intermedio, F1 el peor
\item La diferencia entre F3 y F1 aumenta con riego alto
\item F2 mantiene rendimiento similar independiente del riego
\end{itemize}

\vspace{0.3cm}

\textbf{Implicación práctica:}
\begin{itemize}
\item Si hay riego alto disponible: usar F3
\item Si el riego es limitado: usar F2
\item Evitar F1 salvo restricciones de costo
\end{itemize}
\end{frame}

\subsection{Comparaciones Post-hoc}

\begin{frame}[fragile]{Python: Comparaciones Post-hoc con Pingouin}
\small
\begin{verbatim}
import pingouin as pg

# Crear variable de combinación
df['Combinacion'] = df['Fertilizante'] + '_' + df['Riego']

# Comparaciones múltiples (Tukey HSD)
posthoc = pg.pairwise_tukey(data=df,
                             dv='Rendimiento',
                             between='Combinacion')
print(posthoc[['A', 'B', 'mean(A)', 'mean(B)',
               'diff', 'p-tukey']])
\end{verbatim}

\vspace{0.2cm}
Esto compara todas las combinaciones tratamiento-nivel dos a dos.
\end{frame}

\begin{frame}[fragile]{R: Comparaciones Post-hoc}
\small
\textbf{Con paquete \texttt{emmeans}:}
\begin{verbatim}
library(emmeans)

# Medias marginales estimadas
emm <- emmeans(modelo, ~ fertilizante | riego)
print(emm)

# Comparaciones por pares dentro de cada nivel de riego
pairs(emm)

# O comparar todas las combinaciones
emm_todas <- emmeans(modelo, ~ fertilizante * riego)
pairs(emm_todas, adjust = "tukey")
\end{verbatim}
\end{frame}

\begin{frame}[fragile]{R: Comparaciones con \texttt{agricolae}}
\small
\begin{verbatim}
library(agricolae)

# Crear variable de combinación
df$combinacion <- with(df,
                       interaction(fertilizante, riego))

# ANOVA de un factor con la combinación
modelo_comb <- aov(rendimiento ~ combinacion, data=df)

# Prueba de Tukey
tukey <- HSD.test(modelo_comb, "combinacion",
                  group=TRUE, console=TRUE)

# Mostrar grupos homogéneos
print(tukey$groups)
\end{verbatim}

Esto produce letras que indican grupos que no difieren significativamente.
\end{frame}

\section{Diseño Completamente Aleatorizado (DCA)}

\begin{frame}[fragile]{Ejemplo DCA: Dietas de Pollos}
\small
Un investigador compara 4 dietas en pollos. Peso ganado (gramos) después de 6 semanas:

\begin{verbatim}
Dieta 1: 179, 160, 136, 227, 217, 168, 108, 124
Dieta 2: 309, 229, 181, 141, 260, 203
Dieta 3: 141, 148, 169, 213, 257, 244, 271
Dieta 4: 309, 229, 181, 141, 260, 203, 148
\end{verbatim}

\vspace{0.2cm}
Notemos que hay diferente número de réplicas por tratamiento (diseño no balanceado).
\end{frame}

\begin{frame}[fragile]{Python: ANOVA en DCA}
\small
\begin{verbatim}
import pandas as pd
from scipy import stats

# Crear datos
dieta1 = [179,160,136,227,217,168,108,124]
dieta2 = [309,229,181,141,260,203]
dieta3 = [141,148,169,213,257,244,271]
dieta4 = [309,229,181,141,260,203,148]

# ANOVA de un factor
f_stat, p_valor = stats.f_oneway(dieta1, dieta2,
                                   dieta3, dieta4)
print(f"F = {f_stat:.3f}, p-valor = {p_valor:.4f}")

# Alternativamente, con statsmodels
df = pd.DataFrame({
    'Peso': dieta1 + dieta2 + dieta3 + dieta4,
    'Dieta': ['D1']*8 + ['D2']*6 + ['D3']*7 + ['D4']*7
})
modelo = ols('Peso ~ C(Dieta)', data=df).fit()
print(sm.stats.anova_lm(modelo, typ=2))
\end{verbatim}
\end{frame}

\begin{frame}[fragile]{R: ANOVA en DCA}
\small
\begin{verbatim}
# Crear datos
peso <- c(179,160,136,227,217,168,108,124,
          309,229,181,141,260,203,
          141,148,169,213,257,244,271,
          309,229,181,141,260,203,148)
dieta <- factor(c(rep("D1", 8), rep("D2", 6),
                  rep("D3", 7), rep("D4", 7)))

# ANOVA
modelo_dca <- aov(peso ~ dieta)
summary(modelo_dca)

# Comparaciones múltiples
library(agricolae)
tukey_dca <- HSD.test(modelo_dca, "dieta", group=TRUE)
print(tukey_dca$groups)
\end{verbatim}
\end{frame}

\section{Diseño en Bloques (DBCA)}

\begin{frame}[fragile]{Ejemplo DBCA: Variedades de Trigo}
\small
Rendimiento (ton/ha) de 4 variedades en 5 lotes:

\begin{verbatim}
         Lote1  Lote2  Lote3  Lote4  Lote5
Var1      5.2    4.8    5.5    4.9    5.1
Var2      6.1    5.7    6.4    5.8    6.0
Var3      5.8    5.4    6.1    5.5    5.7
Var4      5.5    5.1    5.8    5.2    5.4
\end{verbatim}
\end{frame}

\begin{frame}[fragile]{Python: ANOVA en DBCA}
\small
\begin{verbatim}
import pandas as pd

# Crear datos
datos_dbca = {
    'Variedad': list('V1')*5 + list('V2')*5 +
                list('V3')*5 + list('V4')*5,
    'Lote': list(range(1,6)) * 4,
    'Rendimiento': [5.2,4.8,5.5,4.9,5.1,
                    6.1,5.7,6.4,5.8,6.0,
                    5.8,5.4,6.1,5.5,5.7,
                    5.5,5.1,5.8,5.2,5.4]
}
df_dbca = pd.DataFrame(datos_dbca)
df_dbca['Lote'] = df_dbca['Lote'].astype('category')

# ANOVA con bloques
modelo_dbca = ols('Rendimiento ~ C(Variedad) + C(Lote)',
                  data=df_dbca).fit()
print(sm.stats.anova_lm(modelo_dbca, typ=2))
\end{verbatim}
\end{frame}

\begin{frame}[fragile]{R: ANOVA en DBCA}
\small
\begin{verbatim}
# Crear datos
rendimiento <- c(5.2,4.8,5.5,4.9,5.1,
                 6.1,5.7,6.4,5.8,6.0,
                 5.8,5.4,6.1,5.5,5.7,
                 5.5,5.1,5.8,5.2,5.4)
variedad <- factor(rep(c("V1","V2","V3","V4"), each=5))
lote <- factor(rep(1:5, 4))

# ANOVA en bloques
modelo_dbca <- aov(rendimiento ~ variedad + lote)
summary(modelo_dbca)

# Comparaciones múltiples (solo variedades)
library(agricolae)
tukey_dbca <- HSD.test(modelo_dbca, "variedad",
                       group=TRUE)
print(tukey_dbca$groups)
\end{verbatim}
\end{frame}

\begin{frame}{Interpretación DBCA}
\textbf{Tabla ANOVA típica:}

\begin{table}
\small
\begin{tabular}{lccccc}
\toprule
\textbf{Fuente} & \textbf{SC} & \textbf{gl} & \textbf{CM} & \textbf{F} & \textbf{p-valor} \\
\midrule
Variedad & 1.275 & 3 & 0.425 & 21.25 & 0.00002 \\
Lote & 0.768 & 4 & 0.192 & 9.60 & 0.00091 \\
Error & 0.240 & 12 & 0.020 & & \\
\midrule
Total & 2.283 & 19 & & & \\
\bottomrule
\end{tabular}
\end{table}

\vspace{0.3cm}

\textbf{Conclusiones:}
\begin{itemize}
\item Hay diferencias significativas entre variedades ($p < 0.001$)
\item El efecto del lote (bloque) también es significativo ($p < 0.001$)
\item El bloqueo fue efectivo (redujo error experimental)
\item Realizar comparaciones post-hoc entre variedades
\end{itemize}
\end{frame}

\section{Diseño en Cuadrado Latino (DCL)}

\begin{frame}[fragile]{Ejemplo DCL: Tiempo de Secado}
\small
Tiempo de secado (minutos) de 4 formulaciones de pintura. Control de operador y día:

\begin{verbatim}
         Lunes  Martes  Mierc.  Jueves
Op1      A:73   B:74    C:71    D:75
Op2      B:73   C:75    D:72    A:74
Op3      C:75   D:73    A:76    B:74
Op4      D:72   A:75    B:73    C:71
\end{verbatim}
\end{frame}

\begin{frame}[fragile]{Python: ANOVA en DCL}
\small
\begin{verbatim}
# Crear datos
datos_dcl = {
    'Tiempo': [73,74,71,75, 73,75,72,74,
               75,73,76,74, 72,75,73,71],
    'Formulacion': list('ABCDBC') + list('DACDABDC'),
    'Operador': [1,1,1,1, 2,2,2,2,
                 3,3,3,3, 4,4,4,4],
    'Dia': [1,2,3,4]*4
}
df_dcl = pd.DataFrame(datos_dcl)
df_dcl['Operador'] = df_dcl['Operador'].astype('category')
df_dcl['Dia'] = df_dcl['Dia'].astype('category')

# ANOVA con tres factores (sin interacción)
modelo_dcl = ols('Tiempo ~ C(Formulacion) + 
                  C(Operador) + C(Dia)',
                  data=df_dcl).fit()
print(sm.stats.anova_lm(modelo_dcl, typ=2))
\end{verbatim}
\end{frame}

\begin{frame}[fragile]{R: ANOVA en DCL}
\small
\begin{verbatim}
# Crear datos
tiempo <- c(73,74,71,75, 73,75,72,74,
            75,73,76,74, 72,75,73,71)
formulacion <- factor(c('A','B','C','D','B','C','D','A',
                        'C','D','A','B','D','A','B','C'))
operador <- factor(rep(1:4, each=4))
dia <- factor(rep(1:4, 4))

# ANOVA en Cuadrado Latino
modelo_dcl <- aov(tiempo ~ formulacion + operador + dia)
summary(modelo_dcl)

# Comparaciones múltiples (solo formulaciones)
library(agricolae)
tukey_dcl <- HSD.test(modelo_dcl, "formulacion",
                      group=TRUE)
print(tukey_dcl$groups)
\end{verbatim}
\end{frame}

\section{Verificación de Supuestos}

\begin{frame}[fragile]{Python: Normalidad}
\small
\begin{verbatim}
from scipy import stats
import matplotlib.pyplot as plt

# Obtener residuos
residuos = modelo.resid

# Prueba de Shapiro-Wilk
stat, p_valor = stats.shapiro(residuos)
print(f"Shapiro-Wilk: W = {stat:.4f}, p = {p_valor:.4f}")

# Gráfico Q-Q
import scipy.stats as stats
fig, ax = plt.subplots(figsize=(8, 6))
stats.probplot(residuos, dist="norm", plot=ax)
ax.set_title("Gráfico Q-Q de Normalidad")
plt.show()

# Histograma
plt.hist(residuos, bins=10, edgecolor='black')
plt.xlabel("Residuos")
plt.ylabel("Frecuencia")
plt.title("Histograma de Residuos")
plt.show()
\end{verbatim}
\end{frame}

\begin{frame}[fragile]{R: Normalidad}
\small
\begin{verbatim}
# Obtener residuos
residuos <- residuals(modelo)

# Prueba de Shapiro-Wilk
shapiro.test(residuos)

# Gráfico Q-Q
qqnorm(residuos)
qqline(residuos, col = "red", lwd = 2)

# Histograma
hist(residuos, breaks = 10,
     main = "Histograma de Residuos",
     xlab = "Residuos", col = "lightblue")
\end{verbatim}

\vspace{0.2cm}
\textbf{Interpretación:}
\begin{itemize}
\item $p > 0.05$ en Shapiro-Wilk: no rechazar normalidad
\item Puntos cercanos a la línea en Q-Q: normalidad razonable
\end{itemize}
\end{frame}

\begin{frame}[fragile]{Python: Homocedasticidad}
\small
\begin{verbatim}
from scipy import stats

# Prueba de Levene
# Agrupar residuos por tratamiento
grupos = [df[df['Fertilizante']=='F1']['Rendimiento'],
          df[df['Fertilizante']=='F2']['Rendimiento'],
          df[df['Fertilizante']=='F3']['Rendimiento']]

stat, p_valor = stats.levene(*grupos)
print(f"Levene: F = {stat:.4f}, p = {p_valor:.4f}")

# Gráfico de residuos vs valores ajustados
valores_ajustados = modelo.fittedvalues
plt.scatter(valores_ajustados, residuos, alpha=0.6)
plt.axhline(y=0, color='red', linestyle='--')
plt.xlabel("Valores Ajustados")
plt.ylabel("Residuos")
plt.title("Residuos vs Valores Ajustados")
plt.show()
\end{verbatim}
\end{frame}

\begin{frame}[fragile]{R: Homocedasticidad}
\small
\begin{verbatim}
library(car)

# Prueba de Levene
leveneTest(modelo)

# Prueba de Bartlett (más sensible a normalidad)
bartlett.test(rendimiento ~ fertilizante)

# Gráfico de residuos vs valores ajustados
plot(modelo, which=1)

# Gráfico de dispersión por grupo
library(ggplot2)
datos <- data.frame(rendimiento, fertilizante)
ggplot(datos, aes(x=fertilizante, y=rendimiento)) +
  geom_boxplot() +
  labs(title="Dispersión por Grupo")
\end{verbatim}

\vspace{0.2cm}
\textbf{Interpretación:}
\begin{itemize}
\item $p > 0.05$ en Levene: no rechazar homocedasticidad
\item Sin patrones en gráfico de residuos: varianza constante
\end{itemize}
\end{frame}

\begin{frame}[fragile]{Python: Independencia}
\small
\begin{verbatim}
from statsmodels.stats.stattools import durbin_watson

# Prueba de Durbin-Watson (autocorrelación)
dw_stat = durbin_watson(residuos)
print(f"Durbin-Watson: {dw_stat:.4f}")

# Gráfico de residuos vs orden
plt.plot(range(len(residuos)), residuos, 'o-')
plt.axhline(y=0, color='red', linestyle='--')
plt.xlabel("Orden de Observación")
plt.ylabel("Residuos")
plt.title("Residuos vs Orden")
plt.show()
\end{verbatim}

\vspace{0.2cm}
\textbf{Interpretación:}
\begin{itemize}
\item DW cercano a 2: no autocorrelación
\item DW cercano a 0: autocorrelación positiva
\item DW cercano a 4: autocorrelación negativa
\end{itemize}
\end{frame}

\begin{frame}[fragile]{R: Independencia}
\small
\begin{verbatim}
library(lmtest)

# Prueba de Durbin-Watson
dwtest(modelo)

# Gráfico de residuos vs orden
plot(residuals(modelo), type='o',
     ylab="Residuos", xlab="Orden de Observación",
     main="Residuos vs Orden")
abline(h=0, col="red", lty=2, lwd=2)
\end{verbatim}

\vspace{0.3cm}
\textbf{Nota importante:}
\begin{itemize}
\item La independencia es un supuesto del diseño experimental
\item La aleatorización adecuada garantiza independencia
\item Las pruebas estadísticas solo detectan dependencia temporal
\item Medidas repetidas o agrupamiento requieren modelos especiales
\end{itemize}
\end{frame}

\section{Transformaciones de Datos}

\begin{frame}{Cuándo y Cómo Transformar}
\textbf{Indicaciones para transformar:}
\begin{itemize}
\item Violación de normalidad o homocedasticidad
\item Relación entre media y varianza de los grupos
\item Mejorar simetría de la distribución
\end{itemize}

\vspace{0.3cm}

\textbf{Transformaciones comunes:}
\begin{table}
\small
\begin{tabular}{ll}
\toprule
\textbf{Situación} & \textbf{Transformación} \\
\midrule
Varianza proporcional a media & Raíz cuadrada: $\sqrt{Y}$ \\
Varianza proporcional a $media^2$ & Logaritmo: $\log(Y)$ \\
Datos de conteo (Poisson) & Raíz cuadrada: $\sqrt{Y + 0.5}$ \\
Proporciones o porcentajes & Arcoseno: $\arcsin(\sqrt{Y})$ \\
Datos muy asimétricos & Box-Cox (óptima) \\
\bottomrule
\end{tabular}
\end{table}
\end{frame}

\begin{frame}[fragile]{Python: Transformación Box-Cox}
\small
\begin{verbatim}
from scipy.stats import boxcox

# Transformación Box-Cox
y_transformado, lambda_opt = boxcox(df['Rendimiento'])
print(f"Lambda óptimo: {lambda_opt:.4f}")

# Crear nueva columna con datos transformados
df['Rendimiento_BC'] = y_transformado

# Re-analizar con datos transformados
modelo_bc = ols('Rendimiento_BC ~ C(Fertilizante) + 
                 C(Riego) + C(Fertilizante):C(Riego)',
                 data=df).fit()
print(sm.stats.anova_lm(modelo_bc, typ=2))

# Verificar supuestos nuevamente
residuos_bc = modelo_bc.resid
stats.shapiro(residuos_bc)
\end{verbatim}
\end{frame}

\begin{frame}[fragile]{R: Transformación Box-Cox}
\small
\begin{verbatim}
library(MASS)

# Encontrar lambda óptimo
bc <- boxcox(modelo, lambda = seq(-2, 2, 0.1))

# Extraer lambda óptimo
lambda <- bc$x[which.max(bc$y)]
print(paste("Lambda óptimo:", round(lambda, 4)))

# Aplicar transformación
if (lambda == 0) {
  y_transformado <- log(rendimiento)
} else {
  y_transformado <- (rendimiento^lambda - 1) / lambda
}

# Re-analizar
modelo_bc <- aov(y_transformado ~ fertilizante * riego)
summary(modelo_bc)

# Verificar supuestos
shapiro.test(residuals(modelo_bc))
\end{verbatim}
\end{frame}

\section{Alternativas No Paramétricas}

\begin{frame}{Pruebas No Paramétricas}
Cuando los supuestos de ANOVA no se cumplen (incluso después de transformaciones):

\vspace{0.3cm}

\begin{table}
\begin{tabular}{ll}
\toprule
\textbf{Diseño} & \textbf{Prueba No Paramétrica} \\
\midrule
DCA (un factor) & Kruskal-Wallis \\
DBCA (bloques) & Friedman \\
DCL (bloques) & Friedman (adaptado) \\
\bottomrule
\end{tabular}
\end{table}

\vspace{0.3cm}

\textbf{Características:}
\begin{itemize}
\item Basadas en rangos, no en valores originales
\item Más robustas ante outliers y no normalidad
\item Menos potentes que ANOVA paramétrica cuando supuestos se cumplen
\item Comparaciones post-hoc con prueba de Dunn o Conover
\end{itemize}
\end{frame}

\begin{frame}[fragile]{Python: Kruskal-Wallis}
\small
\begin{verbatim}
from scipy import stats

# Prueba de Kruskal-Wallis (alternativa a ANOVA de un factor)
dieta1 = [179,160,136,227,217,168,108,124]
dieta2 = [309,229,181,141,260,203]
dieta3 = [141,148,169,213,257,244,271]
dieta4 = [309,229,181,141,260,203,148]

h_stat, p_valor = stats.kruskal(dieta1, dieta2,
                                  dieta3, dieta4)
print(f"Kruskal-Wallis: H = {h_stat:.3f}, 
       p = {p_valor:.4f}")

# Comparaciones post-hoc con pingouin
import pingouin as pg
df_kw = pd.DataFrame({
    'Peso': dieta1 + dieta2 + dieta3 + dieta4,
    'Dieta': ['D1']*8 + ['D2']*6 + ['D3']*7 + ['D4']*7
})
posthoc_kw = pg.pairwise_tests(data=df_kw, dv='Peso',
                                between='Dieta',
                                parametric=False,
                                padjust='bonf')
print(posthoc_kw)
\end{verbatim}
\end{frame}

\begin{frame}[fragile]{R: Kruskal-Wallis y Friedman}
\small
\begin{verbatim}
# Kruskal-Wallis (DCA)
peso <- c(179,160,136,227,217,168,108,124,
          309,229,181,141,260,203,
          141,148,169,213,257,244,271,
          309,229,181,141,260,203,148)
dieta <- factor(c(rep("D1", 8), rep("D2", 6),
                  rep("D3", 7), rep("D4", 7)))

kruskal.test(peso ~ dieta)

# Comparaciones post-hoc
library(FSA)
dunnTest(peso ~ dieta, method="bonferroni")

# Friedman (DBCA)
rendimiento <- matrix(c(5.2,4.8,5.5,4.9,5.1,
                        6.1,5.7,6.4,5.8,6.0,
                        5.8,5.4,6.1,5.5,5.7,
                        5.5,5.1,5.8,5.2,5.4),
                      nrow=4, byrow=TRUE)
friedman.test(rendimiento)
\end{verbatim}
\end{frame}

\section{Resumen y Buenas Prácticas}

\begin{frame}{Flujo de Trabajo Recomendado}
\begin{enumerate}
\item \textbf{Planificación:}
\begin{itemize}
\item Definir hipótesis y diseño experimental
\item Determinar tamaño de muestra
\item Aleatorizar tratamientos
\end{itemize}

\vspace{0.2cm}

\item \textbf{Análisis exploratorio:}
\begin{itemize}
\item Estadísticas descriptivas por grupo
\item Gráficos de caja, medias con IC
\item Identificar outliers
\end{itemize}

\vspace{0.2cm}

\item \textbf{ANOVA:}
\begin{itemize}
\item Ajustar modelo apropiado al diseño
\item Interpretar efectos principales e interacciones
\end{itemize}

\vspace{0.2cm}

\item \textbf{Verificación:}
\begin{itemize}
\item Revisar supuestos (normalidad, homocedasticidad, independencia)
\item Transformar datos si es necesario
\item Considerar alternativas no paramétricas
\end{itemize}
\end{enumerate}
\end{frame}

\begin{frame}{Flujo de Trabajo Recomendado (cont.)}
\begin{enumerate}
\setcounter{enumi}{4}
\item \textbf{Comparaciones post-hoc:}
\begin{itemize}
\item Solo si ANOVA rechaza $H_0$
\item Elegir método apropiado (Tukey, Bonferroni, Scheffé)
\item Ajustar por comparaciones múltiples
\end{itemize}

\vspace{0.2cm}

\item \textbf{Interpretación:}
\begin{itemize}
\item Reportar tamaños de efecto (eta cuadrado, Cohen's d)
\item Graficar resultados con intervalos de confianza
\item Contextualizar conclusiones
\end{itemize}

\vspace{0.2cm}

\item \textbf{Documentación:}
\begin{itemize}
\item Código reproducible (Jupyter, R Markdown)
\item Reportar todos los supuestos verificados
\item Incluir gráficos diagnósticos en apéndices
\end{itemize}
\end{enumerate}
\end{frame}

\begin{frame}{Errores Comunes a Evitar}
\begin{enumerate}
\item \textbf{No verificar supuestos:}
\begin{itemize}
\item Siempre revisar normalidad y homocedasticidad
\item Documentar resultados de pruebas diagnósticas
\end{itemize}

\vspace{0.2cm}

\item \textbf{Ignorar interacciones significativas:}
\begin{itemize}
\item Si hay interacción, no interpretar efectos principales por separado
\item Realizar análisis de efectos simples
\end{itemize}

\vspace{0.2cm}

\item \textbf{No aleatorizar adecuadamente:}
\begin{itemize}
\item La aleatorización es fundamental para validez
\item Documentar proceso de aleatorización
\end{itemize}

\vspace{0.2cm}

\item \textbf{Comparaciones múltiples sin ajuste:}
\begin{itemize}
\item Siempre ajustar p-valores en comparaciones post-hoc
\item Usar métodos apropiados (Tukey, Bonferroni, etc.)
\end{itemize}
\end{enumerate}
\end{frame}

\begin{frame}{Errores Comunes a Evitar (cont.)}
\begin{enumerate}
\setcounter{enumi}{4}
\item \textbf{Confundir significancia estadística con relevancia práctica:}
\begin{itemize}
\item Reportar tamaños de efecto además de p-valores
\item Contextualizar diferencias en términos prácticos
\end{itemize}

\vspace{0.2cm}

\item \textbf{Datos faltantes mal manejados:}
\begin{itemize}
\item En diseños balanceados, datos faltantes complican análisis
\item Considerar imputación o análisis de casos completos
\end{itemize}

\vspace{0.2cm}

\item \textbf{No reportar poder estadístico:}
\begin{itemize}
\item Si no se rechaza $H_0$, considerar poder de la prueba
\item Tamaño de muestra adecuado es fundamental
\end{itemize}

\vspace{0.2cm}

\item \textbf{Uso incorrecto de diseños:}
\begin{itemize}
\item DCL requiere ausencia de interacciones
\item DBCA supone interacción tratamiento-bloque nula
\end{itemize}
\end{enumerate}
\end{frame}

\begin{frame}{Recursos Adicionales}
\textbf{Libros recomendados:}
\begin{itemize}
\item Montgomery, D. C. (2017). \textit{Design and Analysis of Experiments}
\item Kutner et al. (2004). \textit{Applied Linear Statistical Models}
\item Maxwell \& Delaney (2017). \textit{Designing Experiments and Analyzing Data}
\end{itemize}

\vspace{0.3cm}

\textbf{Tutoriales en línea:}
\begin{itemize}
\item Statsmodels documentation (Python)
\item R CRAN Task View: Experimental Design
\item DataCamp, Kaggle, Coursera (cursos de DOE)
\end{itemize}

\vspace{0.3cm}

\textbf{Software especializado:}
\begin{itemize}
\item SAS, SPSS, Minitab (comercial)
\item R packages: agricolae, emmeans, car, lme4
\item Python libraries: statsmodels, pingouin, scikit-posthocs
\end{itemize}
\end{frame}

\begin{frame}{Resumen de la Clase}
\textbf{Implementamos computacionalmente:}
\begin{itemize}
\item ANOVA de dos factores con gráficos de interacción
\item DCA, DBCA, DCL en R y Python
\item Verificación completa de supuestos
\item Comparaciones post-hoc con ajuste por múltiples comparaciones
\item Transformaciones de datos (Box-Cox)
\item Alternativas no paramétricas (Kruskal-Wallis, Friedman)
\end{itemize}

\vspace{0.3cm}

\textbf{Próximos pasos:}
\begin{itemize}
\item Practicar con datasets reales
\item Notebooks de Python y R con ejemplos completos
\item Ejercicios de la Unidad 4
\end{itemize}
\end{frame}

\end{document}
